\chapter{Literature Review}
\label{ch:literature_review}

The rapid expansion of biodiesel production has generated an oversupply of \gls{cgl}, a by-product of the transesterification output, this \Cref{graph:gl-prod} shows the steps to produce biodiesel and when \gls{gl} is producesd \parencite{garlapati2016}. This waste stream, loaded with impurities such as methanol, soaps, and hydroxides, the \Cref{graph:gl-imp} shows the processes and what they contaminate the \gls{gl} with, poses environmental disposal challenges, also with the variation of impurities depending on the source, there's no standard method to use this source of carbon \parencite{anand2012, gebreegziabher2025}. \Glsentrylong{ad} aligns with circular economy principles by converting organic residues into biogas while recycling nutrients, making \gls{cgl} a promising co-substrate \parencite{gebreegziabher2025}.

\begin{figure}[htbp]
    \centering
    \resizebox{0.65\textwidth}{!}{%
        \includestandalone[mode=tex]{assets/graphs/garlapati2016-flow-gl}
    }
    \caption{\glsentrylong{cgl} from biodisel production (Adapted from: \textcite{garlapati2016})}
    \label{graph:gl-prod}
\end{figure}

\begin{figure}[htbp]
    \centering
    \resizebox{\textwidth}{!}{%
        \includestandalone[mode=tex]{assets/graphs/attarbachi2023-flow-gl}
    }
    \caption{Contaminations of the \glsentrylong{cgl} (Adapted from: \textcite{attarbachi2023})}
    \label{graph:gl-imp}
\end{figure}

\textcite{gebreegziabher2025} reported that co-digestion of \gls{cgl} can boost \gls{ch4} production compared to mono-digestion of the base substrate. \textcite{tomczak2025} reviewed one-stage \gls{ad} experiments and concluded that a \gls{ch4} content above \qty{60}{\percent} could be consistently obtained when the \gls{cgl} fraction did not exceed \qty{8}{\percent} (\gls{vpv}) of the feed, the flow chart by \textcite{li2011} on the \Cref{graph:ch4-prod} it shows that using \gls{ad} there could be up to \qty{70}{\percent} \gls{ch4} production. However, the same review emphasized that increasing the \gls{gl} loading beyond a system specific threshold may provoke inhibitory effects, likely due to the accumulation of \glspl{vfa} and the impurities present in \gls{cgl} \parencite{anand2012, tomczak2025}. The heterogeneous composition of \gls{cgl}, affected by oil feedstock, catalyst, and biodiesel processing, demands careful feedstock characterization and, if necessary, simple pretreatment to avoid process upsets \parencite{attarbachi2023}.

\begin{figure}[htbp]
    \centering
    \resizebox{0.65\textwidth}{!}{%
        \includestandalone[mode=tex]{assets/graphs/li2011-flow-ferm}
    }
    \caption{\gls{ch4} production from \gls{ad} (Adapted from: \textcite{li2011})}
    \label{graph:ch4-prod}
\end{figure}

Liquid \gls{ad}, typically conducted in a \gls{cstr}, operates at total solid concentrations of \qtyrange{0.5}{15}{\percent} and is generally used for treating sewage sludge, manure, and \gls{fw} \parencite{li2011}. The \gls{cstr}’s continuous mixing provides uniform distribution of substrates, microorganisms, and heat, which is essential when feeding a fermentable substrate. Temperature is a critical operational parameter; the activity of methanogenic consortia is strongly temperature dependent, and both mesophilic and thermophilic regimes have been applied to \gls{ad} \parencite{choorit2007}. \textcite{tomczak2025} noted that the majority of \gls{gl} co-digestion studies have been performed under mesophilic conditions, while thermophilic operation, which could offer higher conversion rates, still requires more systematic investigation. The \gls{cstr} platform is well suited to explore such temperature effects because its design allows easy control of heating and hydraulic retention time.

The selection of a robust microbial inoculum is decisive for rapid start-up and stable performance of a \gls{cstr} fed with \gls{gl}. Anaerobic sludge from municipal \glspl{wwtp} is the most common source of inoculum because it harbours a diverse consortium of hydrolytic, acidogenic, and methanogenic microorganisms, the flow cahrd from \textcite{gebreegziabher2025} on the \Cref{graph:ch4-steps} shows each step of the \gls{ad} until the \gls{ch4} production. \textcite{posmanik2020} demonstrated that granular sludge, obtained from up-flow anaerobic sludge blanket reactors treating wastewater, outperforms suspended biomass as an inoculum. Although their work focused on batch tests, the findings strongly suggest that using granular sludge from a \gls{wwtp} as seed for a \gls{cstr} could accelerate acclimation, improve process stability, and reduce the risk of acidification during load increases.

\begin{figure}[htbp]
    \centering
    \resizebox{0.65\textwidth}{!}{%
        \includestandalone[mode=tex]{assets/graphs/gebreegziabher2025-flow-ferm}
    }
    \caption{Steps to produce \gls{ch4} from \gls{ad} (Adapted from: \textcite{gebreegziabher2025})}
    \label{graph:ch4-steps}
\end{figure}

When \gls{gl} starts being degraded it makes the \gls{ad} process susceptible to \gls{vfa} accumulation and pH drop if the organic loading rate is not properly controlled. Therefore, close monitoring of process indicators is essential. \textcite{ribas2007} evaluated several methods for determining these parameters specifically in anaerobic reactors, highlighting the importance of accurate \gls{vfa} and \gls{alk} measurements for timely detection of process imbalance. Using the method from DiLallo \& Albertson \parencite{ribas2007} provide a validated method for measuring \glspl{vfa} and \gls{alk}. Maintaining a favourable \gls{vfa_alk} is particularly critical when co-digesting \gls{gl}, as excess \glspl{vfa} can inhibit methanogenesis \parencite{tomczak2025}.

The available literature confirms that \gls{cgl} is a high-energy substrate that can significantly boost \gls{ch4} production in liquid \gls{ad} systems. A \gls{cstr} provides the mixing and temperature control required to harness this potential. Using granular sludge from a \gls{wwtp} as the inoculum can enhance start-up speed and methanogenic activity, as demonstrated by comparative studies. To unlock the full benefits while avoiding inhibition, the \gls{gl} loading must be kept below documented thresholds, and the reactor state should be continuously tracked through standard measurements of \glspl{vfa} and \gls{alk}. Further research is needed to optimise thermophilic \gls{cstr} operation and to establish long-term performance data for granular-sludge-inoculated reactors treating \gls{cgl}.
